\pdfoutput=1

\documentclass[12pt, a4paper]{article}
\usepackage[margin=1.5in, top=1.5in, bottom=1.5in]{geometry}
\usepackage{graphicx}
\usepackage{titlesec}
\usepackage{amsmath}
\usepackage{amssymb}
\usepackage{amsthm}
\usepackage{cases}
\usepackage[english]{babel}
\usepackage{hyperref}
\usepackage{enumitem}
\usepackage{authblk}
\usepackage{dsfont}
\usepackage{sectsty}
\usepackage{xcolor}

\sectionfont{\centering}
% THEOREM / PROPOSITIONS / LEMMAS etc...

\theoremstyle{break}
	\newtheorem{thm}{Theorem}[section]
\theoremstyle{break}
    \newtheorem{prp}{Proposition}[section]
\theoremstyle{break}
	\newtheorem{asm}{Assumption}[section]
\theoremstyle{break}
	\newtheorem{rem}{Remark}[section]
\theoremstyle{break}
	\newtheorem{lem}{Lemma}[section]
\theoremstyle{break}
	\newtheorem{cor}{Corollary}[section]
\theoremstyle{break}
	\newtheorem{defi}{Definition}[section]

% USEFUL COMMANDS

% writing - labels
\makeatletter
\newcommand{\mylabel}[2]{#2\def\@currentlabel{#2}\label{#1}}
\makeatother

\providecommand{\keywords}[1]
{
  \small	
  \textbf{\textit{Keywords}} #1
}

% objects
\newcommand{\Leo}{\mathcal{L}_{\varepsilon}}
\newcommand{\Seq}[2]{\left\{#1_#2\right\}_{#2 = 0}^\infty}

% spaces
\newcommand{\Leb}{L^2\left(\Omega\right)}
\newcommand{\Lei}{L^\infty\left(\Omega\right)}
\newcommand{\CDR}{C^1\left(\mathbb{R}\right)}
\newcommand{\CCSRn}{C_c\left(\mathbb{R}^n\right)}

% operators

\DeclareMathOperator*{\argmax}{arg\,max}
\DeclareMathOperator*{\argmin}{arg\,min}
\DeclareMathOperator{\supp}{supp}
\DeclareMathOperator{\Ima}{Im}
\DeclareMathOperator{\divergence}{div}
\DeclareMathOperator{\Tr}{Tr}

\setlength{\parindent}{0pt}
\counterwithin{equation}{section}

\title{Non-local Klausmeier model of vegetation dynamics in drylands, modeled by bounded, non-periodic regions}

\author[1]{Ricardo Martinez Garcia \thanks{r.martinez-garcia@hzdr.de}}
\author[2]{Michael Hecht \thanks{m.hecht@hzdr.de}}
\author[3]{Maciej Tadej \thanks{maciej.tadej@math.uni.wroc.pl}}

\affil[1,2,3]{\textit{Center for Advanced Systems Understanding (CASUS) – Helmholtz-Zentrum
    Dresden-Rossendorf (HZDR), Untermarkt 20, G\"orlitz 02826, Germany}}
\affil[2,3]{\textit{Insitute of Mathematics, University of Wrocław, pl. Grunwaldzki 2, 50-384 Wrocław, Poland}}

\date{}

\begin{document}

\maketitle

\hspace{20pt}

\begin{abstract}

\end{abstract}

\begin{keywords}
\end{keywords}

\newpage

\section{Introduction}

\subsection{Statement of the problem} We study the following reaction-diffusion-dispersal system of equations
\begin{numcases}{}
    u_t = d_u\mathcal{L} u + u^2 v - B u & \hspace{2mm} $(x, t) \in \Omega \times [0, T],$ \label{eq:u} \\
    v_t = d_v\Delta v - u^2v - v + A & \hspace{2mm} $(x, t) \in \Omega \times [0, T],$ \label{eq:v} \\
    u = 0 & \hspace{2mm} $(x,t) \in \mathbb{R}^n \setminus \Omega  \times [0,T], $ \label{eq:u_bound} \\
    v = 0 & \hspace{2mm} $(x,t) \in \partial\Omega \times [0,T],$ \label{eq:v_bound} \\
    u(0) = u_0 & \hspace{2mm} $x \in \Omega ,$ \label{eq:Unit} \\
    v(0) = v_0 & \hspace{2mm} $x \in \Omega.$ \label{eq:v_init}
\end{numcases}

where the evolution time $T > 0$ and $\Omega \subset \mathbb{R}^n$ is an open and bounded region. The variables $(u, v) := (u(x,t), v(x,t))$ denote the density of a plant and water. Parameters $d_u > 0, d_v > 0$ are dispersal speed and water diffusivity respectively, and parameters $A > 0, B > 0$ are rainfall and plants mortality rates. The dispersal operator is given by the formula
\begin{equation}
    \label{eq:DispersalOperatorDefinition_1}
    \mathcal{L}u(x) = \mathds{1}_{\Omega}(x)\int_{\mathbb{R}^n} J(x-y)\left(u(y) - u(x)\right)\:dy,
\end{equation}
where $\mathds{1}_\Omega$ is an indicator of a domain $\Omega$ and $J : \mathbb{R}^n \rightarrow \mathbb{R}_{\geq 0}$ is an integral kernel falling under the following assumptions.
\begin{asm}
    \label{assumptions_general_kernel}
    We consider the integral kernels $J \in C(\mathbb{R}^n)$ satisfying the following properties
    \begin{enumerate}[label=\textnormal{(\arabic*)}]
        \item [\mylabel{J1}{(J1)}] $J(z) \geq 0$ for $z \in \mathbb{R}^n$ and $J(0) > 0$, 
        \item [\mylabel{J2}{(J2)}] $J(z) = J(-z)$ for $z \in \mathbb{R}^n$,
        \item [\mylabel{J3}{(J3)}] $J(z) \rightarrow 0$ as $|z| \rightarrow \infty$,
        \item [\mylabel{J4}{(J4)}] $\int_{\mathbb{R}^n} J(z)|z|^2dz < \infty$.
        \item [\mylabel{J5}{(J5)}] $\int_{\mathbb{R}^n} J(z) dz = 1$.
    \end{enumerate}
\end{asm}
Observe that assumption \ref{J5} and non-local Dirichlet boundary condition \eqref{eq:u_bound}, when combined, yield the following form of the dispersal operator
\begin{equation}
    \label{EQ:DispersalOperatorDefinition_2}
    \mathcal{L}u(x) = \mathds{1}_{\Omega}(x) \left(\int_{\Omega} J(x-y) u(y)\:dy - u(x)\right).
\end{equation}

\subsection{Ecological background}

\subsection{Related results}

\subsection{Main results}

\section{Preliminaries}
\subsection{Local in-time solutions} 
Here, we construct solutions of the Klausmeier system \eqref{eq:u}-\eqref{eq:v_init}. First, let us consider the space
\begin{equation}
    \label{EQ:C_0_OMEGA}
    C_{0, \Omega}(\mathbb{R}^n) = \left\{ u \in L^\infty(\mathbb{R}^n): u \in C(\Omega), \: u = 0 \text{ in } \mathbb{R}^n \setminus \Omega \right\},
\end{equation}
equipped with usual supremum norm $\|\cdot\|_\infty$.
\begin{lem}
    \label{lem:C_0_OMEGA_BANACH_SPACE}
    $\left(C_{0, \Omega}(\mathbb{R}^n), \|\cdot\|_\infty\right)$ is a Banach space.
\end{lem}
\begin{proof}
We skip the standard proof of the Lemma.
\end{proof}

Now, we extract the linear part in the vegetation equation \eqref{eq:u} and define the corresponding operator $\mathcal{A}: C_{0, \Omega}(\mathbb{R}^n) \mapsto C_{0, \Omega}(\mathbb{R}^n)$ by the formula
\begin{align}
    \label{EQ:A_OPERATOR_DEFINITION}
    \mathcal{A}[u](x) =
    \begin{cases}
        d_u\mathcal{L}u - Bu &x \in \Omega, \\
        0 &x \in \mathbb{R}^n \setminus \Omega.
    \end{cases}
\end{align}
\begin{lem}
    \label{lem:A_BOUNDED_DENSE_LINEAR}
    The operator $\mathcal{A}$ given by formula \eqref{EQ:A_OPERATOR_DEFINITION} is linear and bounded on $C_{0, \Omega}(\mathbb{R}^n)$.
\end{lem}
\begin{proof}
    Using formula \eqref{EQ:DispersalOperatorDefinition_2}, triangle inequality and assumption \ref{J5}, we get that
    \begin{equation*}
        \| \mathcal{A}[u] \|_\infty \leq \| d_u\int_\Omega J(x-y) u(y)\:dy \|_\infty + (B+d_u) \|u\|_\infty \leq (B+2d_u)\|u\|_\infty.
    \end{equation*}
\end{proof}
\begin{lem}
    \label{lem:A_SEMIGROUP_GENERATOR}
    The operator $\mathcal{A}$ given by formula \eqref{EQ:A_OPERATOR_DEFINITION} generates a uniformly continuous semigroup $\{\mathcal{T}_u(t)\}_{t \geq 0}$ of contractions on $C_{0, \Omega}(\mathbb{R}^n)$.
\end{lem}
\begin{proof}
    Lemma \ref{lem:A_BOUNDED_DENSE_LINEAR} yields that $\mathcal{A}$ is bounded, linear operator, thus it generates a uniformly continuous semigroup $\{\mathcal{T}_u(t)\}_{t \geq 0}$. Since $B > 0$ and the spectrum of $\mathcal{L}$ is contained in real left semi-axis, we get that the growth bound $\omega_0$ satisfies
    \begin{equation*}
        \omega_0 := \inf\{\omega \in \mathbb{R}: \lim_{t \rightarrow \infty} e^{-\omega t} \|\mathcal{T}_u(t)\| = 0 \} < 0.
    \end{equation*}
    Hence, $\mathcal{T}_u$ is a contraction semigroup, see Proposition V.1.7 in \cite{EngelNagel2000} for details.
\end{proof}
\begin{cor}
    It immediately follows from Lemma \ref{lem:A_BOUNDED_DENSE_LINEAR}, that $\mathcal{A}$ is an analytic operator.
\end{cor}
Now, we move our attention to the water equation \eqref{eq:v}. Denote by $C_0(\overline{\Omega})$ the Banach space of functions continuous in $\Omega$ and vanishing in $\partial\Omega$. Next, define the operator $\mathcal{B}: D(\mathcal{B}) \mapsto C_0(\overline{\Omega})$ by the formula
\begin{align}
    \label{EQ:B_OPERATOR_DEFINITION}
    \mathcal{B}[v](x) = 
    \begin{cases}
        d_v \Delta v - v &x \in \Omega, \\
        0 &x \in \partial \Omega.
    \end{cases}
\end{align}
To ensure that $\mathcal{B}$ is a closed operator we choose $D(\mathcal{B})$ to be given by
\begin{equation}
    \label{EQ:B_DOMAIN_DEFINITION}
    D(\mathcal{B}) = \left\{ u \in C_0(\overline{\Omega}): \Delta u \in C_0(\overline{\Omega}) \right\}
\end{equation}
\begin{lem}
    \label{lem:B_SEMIGROUP_GENERATOR}
    The operator $\mathcal{B}$ given by formulas \eqref{EQ:B_OPERATOR_DEFINITION} and \eqref{EQ:B_DOMAIN_DEFINITION} generates a strongly continuous semigroup $\{\mathcal{T}_v(t)\}_{t \geq 0}$ of contractions on $C_0(\overline{\Omega})$.
\end{lem}
\begin{proof}
    The proof can be found in the monograph \cite{Arendt2001}. For details, see Theorem 6.1.8, p. 400.
\end{proof}
Now, we are ready to construct the solutions of the Klausmeier system \eqref{eq:u}-\eqref{eq:v_init}.
\begin{thm}[Local existence]
    \label{thm:Well-Posedness}
    Let $(u_0, v_0) \in C_{0, \Omega}(\mathbb{R}^n) \times C_0(\overline{\Omega})$. Then there exists $T > 0$ and a unique solution $(u, v)$ to the problem \eqref{eq:u}-\eqref{eq:v_init}, such that
    \begin{equation}
        u \in C\left([0, T], C_{0, \Omega}(\mathbb{R}^n)\right) \cap C^1\left([0, T), C_{0, \Omega}(\mathbb{R}^n)\right),
    \end{equation}
    and
    \begin{equation}
        v \in C\left([0, T], C_0(\overline{\Omega})\right) \cap C^1\left([0, T), C_0(\overline{\Omega})\right).
    \end{equation}
\end{thm}
\begin{proof}
    Given $\mathcal{A}, \mathcal{B}$ by the formulas \eqref{EQ:A_OPERATOR_DEFINITION} and \eqref{EQ:B_OPERATOR_DEFINITION}-\eqref{EQ:B_DOMAIN_DEFINITION} respectively, let us define the operator $\mathcal{C}: C_{0, \Omega}(\mathbb{R}^n) \times D(\mathcal{B}) \mapsto C_{0, \Omega}(\mathbb{R}^n) \times C_0(\overline{\Omega})$ by the formula
    \begin{equation}
        \mathcal{C} = 
        \begin{bmatrix}
            \mathcal{A} & 0\\ 0 & \mathcal{B}
        \end{bmatrix}.
    \end{equation}
    Next, denote the vectors $w = (u, v)$ and $w_0 = (u_0, v_0)$. Using this notation we re-write the Klausmeier system \eqref{eq:u}-\eqref{eq:v_init} into
    \begin{equation}
        \label{eq:ClassicalSemigroupFormulation}
         w'(t) = \mathcal{C} w(t) + F(w(t)), \quad w(0) = w_0.
    \end{equation}
    where nonlinearity $F: C_{0, \Omega}(\mathbb{R}^n) \times C_0(\overline{\Omega}) \mapsto C_{0, \Omega}(\mathbb{R}^n) \times C_0(\overline{\Omega})$ is given by
    \begin{align}
        F(w) = F(u, v) = 
        \begin{cases}
            (u^2 v, -u^2 v + A) &x \in \Omega, \\
            (0, 0) &x \notin \Omega
        \end{cases}.
    \end{align}
    Lemma \ref{lem:A_SEMIGROUP_GENERATOR} and Lemma \ref{lem:B_SEMIGROUP_GENERATOR} give that $\mathcal{C}$ generates a semigroup $\{\mathcal{T}(t)\}_{t \geq 0}$ of contractions on $C_{0, \Omega}(\mathbb{R}^n) \times C_0(\overline{\Omega})$. Thus, we construct the solution $w$ from the integral equation
    \begin{equation}
        \label{eq:MildSemigroupFormulation}
         w(t) = \mathcal{T}(t) w_0 + \int_0^t \mathcal{T}(t-s) F(w(s)) \:ds.
    \end{equation}
    Since $w_0 \in C_{0, \Omega}(\mathbb{R}^n) \times C_0(\overline{\Omega})$ and $F$ is a locally Lipschitz function, we proceed with a standard fixed point argument to obtain that there exists positive number $T := T(w_0)$ such that \eqref{eq:MildSemigroupFormulation} admits a unique solution $w \in C\left([0, T], C_{0, \Omega}(\mathbb{R}^n) \times C_0(\overline{\Omega})\right)$. To show differentiability in time one needs to integrate equations \eqref{eq:u}, \eqref{eq:v} to find that $w$ satisfies
    \begin{equation}
        \label{eq:IntegralFormula}
        w(t) = w_0 + \int_0^t \mathcal{C}w(s) + F(w(s)) \:ds.
    \end{equation}
    It is an immediate consequence of the continuity of $w$, that the formula \eqref{eq:IntegralFormula} yields differentiability of $w$ over time interval $[0, T)$.
\end{proof}
\subsection{Non-negativity and bounds} 
Here, we are interested in non-negativity of solutions obtained with Theorem~\ref{thm:Well-Posedness} and existence of an invariant region $\Gamma$ for the system \eqref{eq:u}-\eqref{eq:v_init}.
\begin{thm}[Non-negativity and invariant region]
    \label{thm:Non-negativity}
     If the initial datum $(u_0, v_0) \in C_{0, \Omega}(\mathbb{R}^n) \times C_0(\overline{\Omega})$ is non-negative then the solution $(u, v)$ constructed in Theorem \ref{thm:Well-Posedness} is non-negative and for all $(x, t) \in \Omega \times [0, T]$ it is true that
     \begin{equation}
         v(x,t) \leq \max\{\|v_0\|_\infty, A\} =: R_1.
     \end{equation}
     Moreover, denote the set $\Gamma = [0, B/R_1]$. If $u_0 \in \Gamma$ for $x \in \Omega$ then $u \in \Gamma$ for all $(x, t) \in \Omega \times [0, T]$.
\end{thm}
\begin{proof}
    It follows from the parabolic maximum principle that for all $(x, t) \in \Omega \times [0, T)$ we have the estimate of the form
    \begin{equation}
        \label{EQ:AprioriEstimate_V}
        0 \leq v(x, t) \leq \max\{\|v_0\|_\infty, A\} =: R_1.
    \end{equation}
    Next, we define the point $\xi := \xi(t) \in \Omega$ and evaluation $\mu := \mu(t)$ such that 
    \begin{equation*}
        \mu(t) := u(\xi(t), t) = \min_{x \in \Omega} u(x, t).
    \end{equation*}
    Without loss of generality, we take arbitrary $0 < s < t < T$ and estimate
    \begin{equation*}
        \left|\mu(t) - \mu(s)\right| \leq \left|u(\xi(s), t) - u(\xi(s), s)\right| \leq \max_{(x, r) \in \Omega \times [0, T)} |u(x, r)| |t-s|.
    \end{equation*}
    This is Lipschitz continuity of $\mu$. It follows from the Radamacher Theorem that $\mu$ is differentiable almost everywhere in $[0, T]$. Next, we estimate the difference quotients applying Mean Value Theorem
    \begin{align*}
        \frac{\mu(t+h) - \mu(t)}{h} &\leq \frac{u(\xi(t+h), t+h) - u(\xi(t + h), t)}{h} = u_t(\xi(t + h), c(t, h)), \\
        \frac{\mu(t+h) - \mu(t)}{h} &\geq \frac{u(\xi(t), t+h) - u(\xi(t), t)}{h} = u_t(\xi(t), c(t, h)),
    \end{align*}
    where $c(t,h) \in [t, t+h]$ is an intermediate point. Passing with $h \rightarrow 0$ and estimating $\mathcal{L}$ in a minimum point, we obtain the differential inequality
    \begin{align}
        \label{eq:differential_inequality_min}
        \begin{split}
            \mu'(t) &= \mathcal{L} u(\xi(t), t) + \mu^2(t) v(\xi(t), t) - B \mu(t) \geq - B \mu(t), \\
            \mu(0) &= \min_{x \in \Omega} u_0(x) \geq 0.
        \end{split}
    \end{align}
    When solved, inequality \eqref{eq:differential_inequality_min} yields the following lower bound 
    \begin{equation*}
        u(x,t) \geq \mu(t) \geq \mu(0) e^{-B t} \geq 0.
    \end{equation*}
    Now, we define the point $\zeta := \zeta(t) \in \Omega$ and evaluation $\nu := \nu(t)$ such that
    \begin{equation*}
        \nu(t) := u(\zeta(t), t) = \max_{x \in \Omega} u(x,t).
    \end{equation*}
    Proceeding as previously, we arrive at the following differential inequality
    \begin{align}
        \label{eq:differential_inequality_max}
        \begin{split}
            \nu'(t) &= \mathcal{L} u(\zeta(t), t) + \nu^2(t) v(\zeta(t), t) - B \nu(t) \leq R_1 \nu^2(t) - B \nu(t), \\
            \nu(0) &= \max_{x \in \Omega} u_0(x) \leq B/R_1.
        \end{split}
    \end{align}
    When solved, inequality \eqref{eq:differential_inequality_max} yields the following upper bound
    \begin{equation}
        \label{EQ:AprioriEsitmate_Unvariant_Region}
        u(x,t) \leq \nu(t) \leq \frac{B\nu(0)}{R_1 \nu(0) + 
        (B - R_1 \nu(0)) e^{B t}} < B / R_1.
    \end{equation}
    Hence $\Gamma = [0, B/R_1]$ is an invariant region.
\end{proof}

\section{Stationary solutions}

We consider the stationary solutions of \eqref{eq:u}-\eqref{eq:v_init} with a corresponding system of equations
\begin{numcases}{}
    0 = d_u\mathcal{L} u + u^2 v - B u & \hspace{2mm} $x \in \Omega,$ \label{eq:u_stationary} \\
    0 = d_v\Delta v - u^2v - v + A & \hspace{2mm} $x \in \Omega,$ \label{eq:v_stationary} \\
    u = 0 & \hspace{2mm} $x \in \mathbb{R}^n \setminus \Omega , $ \label{eq:u_bound_stationary} \\
    v = 0 & \hspace{2mm} $x \in \partial\Omega,$ \label{eq:v_bound_stationary}
\end{numcases}

\subsection{Reduction of the problem} Our first goal is to reduce the number of equations in the system \eqref{eq:u_stationary}-\eqref{eq:v_bound_stationary}.
\begin{lem}
    \label{lem:existence_of_inhomogeneous_elliptic_v}
    Let $u \in L^\infty(\Omega)$ be a given function. Then, there exists a unique, weak solution $V \in W^{1,2}_0(\Omega) \cap W^{2,2}(\Omega)$ of equations \eqref{eq:v_stationary} and \eqref{eq:v_bound_stationary}. Moreover, we have that $\|V \|_\infty \leq A$.
\end{lem}
\begin{proof}
    Consider the following weak formulation of equations \eqref{eq:v_stationary}-\eqref{eq:v_bound_stationary}
    \begin{equation}
        \label{eq:weak_formulation_v}
        d_v\int_\Omega\nabla v \cdot \nabla w \:dx + \int_\Omega (u^2 + 1) v w\:dx = A \int_\Omega w \:dx,
    \end{equation}
    where $w \in W^{1,2}_0(\Omega)$ is arbitrary. We immediately estimate the left hand side of equation \eqref{eq:weak_formulation_v} as follows
    \begin{equation}
        \label{eq:coercivity_bilinear_form_v}
        d_v\int_\Omega | \nabla v |^2 \:dx + \int_\Omega (u^2 + 1) v^2\:dx \geq \min\{d_v, 1\} \|v\|_{1,2}^2
    \end{equation}
    Using Cauchy-Schwarz and Poincare inequalities, one bounds the right hand side of the equation \eqref{eq:weak_formulation_v} by the following expression
    \begin{equation}
        \label{eq:boundedness_bilinear_form_v}
        A \int_\Omega v \:dx \leq A |\Omega|^{1/2} \|v\|_2 \leq A |\Omega|^{1/2} \|v\|_{1,2},
    \end{equation}
    This shows that the left hand side of the equation \eqref{eq:weak_formulation_v} defines a bounded and coercive bilinear form. From the Lax-Milgram Theorem we have that there exists a unique solution $V \in W^{1,2}_0(\Omega)$. 

    Now, we argue that in fact there is more regularity. Indeed let us re-write the equation \eqref{eq:v_stationary} as follows
    \begin{equation}
        \label{eq:laplacian_stationary}
        \Delta v (x) = (u^2(x) + 1) v(x)  - A, \quad x \in \Omega.
    \end{equation}
    we have already shown that $v \in W^{1,2}_0(\Omega)$. Hence, the expression $(u^2 + 1)v - A \in L^2(\Omega)$ and consequently, from equation \eqref{eq:laplacian_stationary}, we get that $\Delta v \in L^2(\Omega)$. This yields that $v \in W^{2,2}(\Omega)$ as well.
    
    Finally, it is worth to notice that the upper bound $\|V\|_\infty \leq A$ follows directly from the elliptic maximum principle.
\end{proof}

\begin{lem}
    \label{eq:stationary_v_lipschitz_constant}
    Consider $V : B_R \rightarrow W^{1,2}_0(\Omega)$, where $B_R \subset L^\infty(\Omega)$ is a closed ball with radius $R > 0$. We define $V$ to be a weak solution of the problem \eqref{eq:weak_formulation_v}. It is a Lipchitz function, with a Lipschitz constant given by
    \begin{equation}
        Lip(V) = \frac{2R C A |\Omega|^{1/2}}{\min\{d_v, 1\}}
    \end{equation}
\end{lem}

\begin{proof}
    First, denote $u, v \in L^\infty(\Omega)$. We reformulate the weak problem \eqref{eq:weak_formulation_v} in terms of the difference $W = V(u) - V(v)$. Let $w \in W^{1,2}_0(\Omega)$, we write
    \begin{equation*}
        d_v \int_\Omega \nabla W \cdot \nabla w \:dx + \int_\Omega (u^2 + 1) V(u) w \:dx - \int_\Omega (v^2 + 1) V(v) w \:dx = 0
    \end{equation*}
    This, may be re-written, by adding $0 = V(v) - V(v)$ to the second term
    \begin{align}
        \label{eq:lipschitz_continuity_weak_formulation_difference}
        d_v \int_\Omega \nabla W \cdot \nabla w \:dx &+ \int_\Omega (u^2 + 1) W w \:dx = \int_\Omega V(v) (u^2 - v^2) w \:dx.
    \end{align}
    Plug $w = W$, we estimate the left hand side of \eqref{eq:lipschitz_continuity_weak_formulation_difference} from below
    \begin{align}
        \label{eq:ineaulity_lhs_lipschitz_cont}
        \begin{split}
            d_v \int_\Omega |\nabla W|^2\:dx + \int_\Omega (u^2 + 1) W^2 \:dx
            &\geq
            \min\{d_v, 1\} \| W_M \|_{1, 2}^2.
        \end{split}
    \end{align}
    Now, we apply the Cauchy-Schwarz and Poincare inequalities together with Lemma \ref{lem:existence_of_inhomogeneous_elliptic_v} applied to the right hand side of equation \eqref{eq:lipschitz_continuity_weak_formulation_difference}, to obtain
    \begin{align}
        \label{eq:ineaulity_rhs_lipschitz_cont}
        \begin{split}
            \int_\Omega V(v) \left( u^2 - v^2 \right) W  \:dx 
            &\leq
            \|V(v)\|_\infty (\|u\|_\infty + \|v\|_\infty) \int_\Omega W\:dx \|u-v\|_\infty\\
            &\leq
            2RA|\Omega|^{1/2}\|W\|_{1,2}\|u-v\|_\infty 
        \end{split}
    \end{align}
    Applying the Sobolev embedding $W^{1,2}_0(\Omega) \hookrightarrow L^\infty(\Omega)$ to the combined inequalities \eqref{eq:ineaulity_lhs_lipschitz_cont} and \eqref{eq:ineaulity_rhs_lipschitz_cont} we get that
    \begin{equation*}
        \|W\|_\infty = \|V(u) - V(v) \|_\infty \leq \frac{2RCA|\Omega|^{1/2}}{\min\{d_v, 1\}}\|u-v\|_\infty,
    \end{equation*}
    where $C$ is a constant emerging from the Sobolev embedding.
\end{proof}

\subsection{Existence of solutions} Now, we reduce the problem \eqref{eq:u_stationary}-\eqref{eq:v_bound_stationary} into the following form
\begin{numcases}{}
    0 = d_u\mathcal{L} u + u^2 V(u) - B u & \hspace{2mm} $x \in \Omega,$ \label{eq:u_stationary_reduced} \\
    u = 0 & \hspace{2mm} $x \in \mathbb{R}^n \setminus \Omega , $ \label{eq:u_bound_stationary_reduced}
\end{numcases}
where we denoted by $V := V(u)$ the solution obtained in Lemma \ref{lem:existence_of_inhomogeneous_elliptic_v}. We will construct solution $(u, V(u))$ of the reduced system \eqref{eq:u_stationary_reduced}-\eqref{eq:u_bound_stationary_reduced} in the neighborhood of constant solutions of the following system of equations.
\begin{align}
    \label{EQ:ConstantSteadyStates}
    \begin{cases}
        0 = u^2v - Bu,\\
        0 = -u^2v - v + A.
    \end{cases}
\end{align}
In particular these are the constant solutions of equations \eqref{eq:u_stationary} and \eqref{eq:v_stationary}. We start with the following fact.
\begin{lem}
    \label{lem:stationary_v_estimate_v_star}
    Let $(u_*, v_*)$ be a real solution of the system \eqref{EQ:ConstantSteadyStates}. Denote by $V(u_*)$ the solution of equations \eqref{eq:v_stationary} and \eqref{eq:v_bound_stationary}. For all $\epsilon > 0$ there exists diffusion rate $d_v > 0$, sufficiently small, such that
    \begin{equation*}
        \| V(u_*) - v_*\|_2 \leq \epsilon.
    \end{equation*}
\end{lem}

\begin{proof}
    First, we denote the orthonormal set $\{\phi_i\}_{i \in \mathbb{N}}$ in $L^2(\Omega)$. We choose it to be the set of the eigenfunctions of $\Delta$ with Dirichlet boundary condition. Next, 
    we expand a constant function $A = \sum_{j=1}^\infty a_j \phi_j$. Since $V \in L^2(\Omega)$ we can expand it in the eigenfunction basis as well. Namely $V(u_*) = \sum_{j=1}^\infty c_j \phi_j$. Plugging this into equation \eqref{eq:v_stationary} we get that
    \begin{equation*}
        d_v\sum_{i=1}^\infty c_j \lambda_j \phi_j - (u_*^2 + 1) \sum_{j=1}^\infty c_j \phi_j + \sum_{j=1}^\infty a_j \phi_j = 0.
    \end{equation*}
    From the orthogonality, we get that the following equation holds for $j \in \mathbb{N}$
    \begin{equation*}
        d_v c_j\lambda_j - (u_*^2+1) c_j + a_j = 0,
    \end{equation*}
    solving for coefficients $c_j$, we obtain the formula
    \begin{equation*}
        c_j = \frac{a_j}{u_*^2 + 1 - d_v\lambda_j}
    \end{equation*}
    It follows that
    \begin{equation*}
        V(u_*) = \sum_{j=1}^\infty \frac{a_j}{u_*^2 + 1 - d_v\lambda_j} \phi_j.
    \end{equation*}
    Next, we plug the eigenfunction expansion of a constant function $A$ to the definition of steady state $v_* = \frac{A}{u_*^2 + 1}$ to obtain the following estimate
    \begin{align*}
        \|V(u_*) - v_*\|_2 
        &=
        \|\sum_{j=1}^\infty \frac{a_j}{u_*^2 + 1 - d_v\lambda_j} \phi_j - \sum_{j=1}^\infty \frac{a_j}{u_*^2 + 1}\phi_j\|_2  \\
        &=
        \|\sum_{j=1}^\infty\left(\frac{1}{u_*^2+1} - \frac{1}{u_*^2+1-d_v\lambda_j} \right) a_j\phi_j\|_2  \\
        &\leq
        \sum_{j=1}^\infty\left| \frac{d_v \lambda_j a_j}{(u_*^2+1)(d_v\lambda_j - (u_*^2+1))} \right| \\
        &=
        \sum_{j=1}^\infty\left| \frac{a_j}{(u_*^2+1)(1 - (u_*^2+1) / d_v \lambda_j)} \right|.
    \end{align*}
    For any $\epsilon > 0$ there exists $N \in \mathbb{N}$ such that $\sum_{j=N+1}^\infty |a_j| < \epsilon$. Next, we take $d_v < \frac{C_N}{N \lambda_N} \epsilon$ where constant $C_N >0$ will be specified later on. We estimate as follows
    \begin{align*}
        \|V(u_*) - v_*\|_2
        &\leq
        \sum_{j=1}^N\left| \frac{d_v \lambda_j a_j}{(u_*^2+1)(d_v\lambda_j - (u_*^2+1))} \right| \\
        &+
        \sum_{j=N+1}^\infty\left| \frac{a_j}{(u_*^2+1)(1 - (u_*^2+1) / d_v \lambda_j)} \right| \\
        &\leq
        C_N \epsilon \left| \frac{a_{j'}}{(u_*^2+1)(d_v\lambda_{j'} - (u_*^2+1))} \right| \\
        &+
        \frac{C_N}{N} \epsilon \left| \frac{1}{(u_*^2+1)(1 - u_*^2+1) / d_v \lambda_{N+1})} \right| \\
        &= \epsilon.
    \end{align*}
    where $j' \in \{1, \ldots, N\}$ is such that $\|\{a_j\}_{j \in \mathbb{N}}\|_\infty = a_{j'}$. Note that in order to have that $j' \in \{1, \ldots, N\}$, we need to take $\epsilon$ sufficiently small. Constant $C_N > 0$ is given explictly by
    \begin{equation*}
        C_N = \left( \left| \frac{a_{j'}(u_*^2+1)^{-1}}{((d_v\lambda_{j'} - (u_*^2+1))} \right| + \left| \frac{(u_*^2+1)^{-1}}{N((1 - u_*^2+1) / d_v \lambda_{N+1})} \right| \right)^{-1}.
    \end{equation*}
\end{proof}

\begin{thm}
    \label{thm:Stationary-Solutions-Existence}
    Denote by $(u_*, v_*)$ a solution of the system \eqref{EQ:ConstantSteadyStates}. For every parameters $A > 0, B > 0$ there is a dispersal rate $d_u >0$ large, such that there exists a unique solution $(U, V)$ of equations \eqref{eq:u_stationary}-\eqref{eq:v_bound_stationary} satisfying
    \begin{equation}
        \|U - u_*\|_{L^\infty(\Omega)} \leq R, \quad \|V - v_*\|_{L^2(\Omega)} \leq Lip(V) |\Omega|^{1/2} R + \epsilon,
    \end{equation}
    for $R > 0$ sufficiently small and some $\epsilon := \epsilon(d_v) > 0$.
\end{thm}

\begin{proof}
    Let $U:= U(x)$ be the discontinuous function defined by the formula
    \begin{align}
        \label{eq:stationary_u_definition}
        U(x) 
        = 
        u_* \mathds{1}_{\Omega}(x) + \varphi(x)
    \end{align}
    where $\varphi \in C_{0, \Omega}(\mathbb{R}^n)$. We plug that function into equation \eqref{eq:u_stationary} to get  
    \begin{equation}
        \label{eq:stationary_u_plugged}
        0 = d_u \mathcal{L} U + U^2 V(U) - B U.
    \end{equation}
    We take $x \in \Omega$ and re-write the non-local term in \eqref{eq:stationary_u_plugged} using property \ref{J5} and equivalence of definitions \eqref{eq:DispersalOperatorDefinition_1}, \eqref{EQ:DispersalOperatorDefinition_2} of the dispersal operator.
    \begin{align}
        \label{eq:non-local-part-u-reformulation}
        \begin{split}
        \mathcal{L}U 
        &=
            \mathcal{L}\varphi + u_*  \int_{\mathbb{R}^n} J(x-y)\left(\mathds{1}_{\Omega}(y) - \mathds{1}_{\Omega}(x)\right)\:dy  \\
            &= \mathcal{L}\varphi + u_* \left( \int_\Omega J(x-y)\:dy - 1 \right) \\
            &= \mathcal{L}\varphi - u_*  \int_{\mathbb{R}^n \setminus \Omega} J(x-y)\:dy =: \mathcal{L}\varphi - u_* j_\Omega(x)
        \end{split}
    \end{align}
    Finally, applying formula \eqref{EQ:DispersalOperatorDefinition_2} and assumption $\varphi \in C_{0, \Omega}(\mathbb{R}^n)$ to the first term, we obtain that
    \begin{equation}
        \label{eq:non-local-part-u-stationary-plugged}
        \mathcal{L}\varphi =  \int_{\Omega} J(x-y) \varphi(y)\:dy - \varphi(x) =: J *\varphi(x) - \varphi(x)
    \end{equation}
    Since $-u_*^2v_* + Bu_*=0$, we re-write the non-linear part of equation \eqref{eq:stationary_u_plugged}
    \begin{align}
        \label{eq:non-linear-part-u-stationary-plugged}
        \begin{split}
            U^2 V(U) - B U 
            &= U^2 V(U) - u_*^2 v_* - B (U - u_*) \\
            &= U^2 V(U) - u_*^2 V(U) + u_*^2 V(U) - u_*^2 v_* - B \varphi  \\
            &= (U^2 - u_*^2)V(U) - u_*^2(V(U) - v_*) - B \varphi \\
            &= (2 u_* \varphi + \varphi^2) V(U) - u_*^2(V(U) - v_*) - B \varphi.
        \end{split}
    \end{align}
    Combining formulas \eqref{eq:non-local-part-u-reformulation}, \eqref{eq:non-local-part-u-stationary-plugged} and \eqref{eq:non-linear-part-u-stationary-plugged}, we re-formulate the stationary problem \eqref{eq:stationary_u_plugged} into the following fixed point problem.
    \begin{equation}
        \label{eq:u-stationary-fixed-point}
        \varphi = \frac{1}{d_u + B}\Big(d_u J *\varphi - d_u u_* j_\Omega + (2 u_* \varphi + \varphi^2) V(U) - u_*^2(V(U) - v_*)\Big)
    \end{equation}
    Denote the Banach space $X = \left\{ \varphi \in C_{0,\Omega}(\mathbb{R}^n): \|\varphi\|_\infty \leq R \right\}$. Then, define an operator $\mathcal{F}: X \rightarrow X$, corresponding to equation \eqref{eq:u-stationary-fixed-point}, by equation
    \begin{align}
        \label{eq:operator-f-def}
        \begin{split}
            \mathcal{F}(\varphi) 
            &= 
            \frac{d_u \mathds{1}_{\Omega}}{d_u + B}\left( J *\varphi - u_* j_\Omega\right) + \frac{\mathds{1}_{\Omega}}{d_u + B}(2 u_* \varphi + \varphi^2) V (u_* + \varphi)\\
            &- \frac{u_*^2\mathds{1}_{\Omega}}{d_u + B}(V(u_* + \varphi) - v_*)        
        \end{split}
    \end{align}
    First, we verify that indeed $\mathcal{F}: X \rightarrow X$. Let $\varphi \in X$, we have that
    \begin{align}
        \label{eq:f_ball_to_ball_condition}
        \begin{split}
            \|\mathcal{F}(\varphi)\|_{\infty}
            &\leq
            \frac{d_u }{d_u + B} \left( \|1 -j_\Omega\|_\infty R + u_*\|j_\Omega\|_\infty\right) \\
            &+
            \frac{A}{d_u + B}\left( 2 u_* R + R^2 \right) 
            + \frac{u_*^2 (A + v_*)}{d_u + B} \leq R.
        \end{split}
    \end{align}
    This condition holds whenever the dispersal rate $d_u > 0$ is large enough. Having that, we show the contractivity of $\mathcal{F}$. Suppose that $\varphi, \psi \in X$, then
    \begin{align}
        \label{eq:contraactivity-first-inequality}
        \begin{split}
            \|\mathcal{F}(\varphi) - \mathcal{F}(\psi)\|_\infty
            &\leq
            \frac{d_u}{d_u + B} \|1 - j_\Omega\|_\infty \|\varphi - \psi\|_\infty \\
            &+ \frac{1}{d_u + B} \left\|(2u_* \varphi + \varphi^2) V(u_* + \varphi)
            - (2u_* \psi + \psi^2) V(u_* + \psi) \right\|_\infty \\
            &+ \frac{u_*^2}{d_u + B}\left\| V(u_* + \varphi) - V(u_* + \psi) \right\|_\infty.
        \end{split}
    \end{align}
    Consider the function $g(\varphi) = (2u_*\varphi + \varphi^2)V(u_* + \varphi)$. Applying the Mean Value Theorem and Lemmas \ref{eq:stationary_v_lipschitz_constant}, \ref{lem:stationary_v_estimate_v_star} we obtain the following bound
    \begin{align*}
        \|g'(\varphi)\|_\infty 
        &= 
        \|(2u_* + 2 \varphi)V(u_* + \varphi) + (2u_* \varphi + \varphi^2) \frac{\partial V}{\partial \varphi}(u_* + \varphi)\|_\infty \\
        &\leq
        (2u_* + 2R) A + (2u_*R + R^2) Lip(V).
    \end{align*}
    This yields that 
    \begin{equation}
        \label{eq:contractivity-second-term}
        \|g(\varphi) - g(\psi)\|_\infty \leq \left( (2u_* + 2R) A + (2u_*R + R^2) Lip(V) \right) \|\varphi - \psi\|_\infty.
    \end{equation}
    The last term in inequality \eqref{eq:contraactivity-first-inequality} is estimated using Lemma \ref{lem:stationary_v_estimate_v_star}, namely
    \begin{equation}
        \label{eq:contracitivity-third-term}
        \left\| V(u_* + \varphi) - V(u_* + \psi) \right\|_\infty
        \leq
        Lip(V)\|\varphi - \psi\|_\infty
    \end{equation}
    Now, combining inequality \eqref{eq:contraactivity-first-inequality} with \eqref{eq:contractivity-second-term} and \eqref{eq:contracitivity-third-term} we obtain
    \begin{align}
        \label{eq:contractivity-final-inequality}
        \begin{split}
            \|\mathcal{F}(\varphi) - \mathcal{F}(\psi)\|_\infty
            &\leq
            \frac{d_u}{d_u + B} \|1 - j_\Omega\|_\infty \|\varphi - \psi\|_\infty\\
            &+
            \frac{(2u_* + 2R) A + (2u_*R + R^2) Lip(V)}{d_u + B} \|\varphi - \psi\|_\infty \\
            &+ \frac{u_*^2Lip(V)}{d_u + B}\|\varphi - \psi\|_\infty.
        \end{split}
    \end{align}
    Provided that $d_u > 0$ is large enough and $R > 0$ is small enough we get contractivity of $\mathcal{F}$. Now, apply Banach Fixed Point Theorem to obtain existence of a unique solution $(U, V)$. From the construction it is clear that 
    \begin{equation}
        \max_{x \in \Omega} |U(x) - u_*| \leq R.
    \end{equation}
    Moreover, from the Lemma \ref{lem:stationary_v_estimate_v_star} it immediately follows that
    \begin{align}
        \begin{split}
            \|V(U) - v_*\|_2 
            &\leq \|V(U) - V(u_*)\|_2 + \|V(u_*) - v_*\|_2 \\
            &\leq Lip(V)\|U-u_*\|_2 + \epsilon \\
            &\leq Lip(V)|\Omega|^{1/2} R + \epsilon.     
        \end{split}
    \end{align}
\end{proof}

\begin{rem}
    From Lemma \ref{lem:stationary_v_estimate_v_star} and Theorem \ref{thm:Stationary-Solutions-Existence}, it follows that small $d_v$ yields existence of the water profile $V$, which is a small perturbation of the constant number $v_*$. However, ecologically, vegetation survival is more critical, which occurs regardless of the magnitude of $d_v$. 
\end{rem}

\section{Extinction and survival of vegetation}

\subsection{Insufficient initial biomass} Here, we investigate how the assumption of small initial density leads to the disappearance of vegetation density and stabilization of water density.

\begin{thm}
    Let $R_1 = \max\{\|v_0\|_\infty, A\}$. Denote an interval $\Gamma  = \left[0, \frac{B}{R_1}\right]$. Suppose that $u_0 \in \Gamma$ for all $x \in \Omega$. Then the solution $(u, v)$ of the problem \eqref{eq:u}-\eqref{eq:v_init} satisfies
    \begin{equation}
        u(x,t) \xrightarrow{t \rightarrow \infty} 0, \quad v(x,t) \xrightarrow{t \rightarrow \infty} V_0(x),
    \end{equation}
    where we denoted the water profile, obtained in Lemma \ref{lem:existence_of_inhomogeneous_elliptic_v}, by $V_0 := V(0)$.
\end{thm}
\begin{proof}
    The first limit is the consequence of the bound \eqref{EQ:AprioriEsitmate_Unvariant_Region}. Since $(B - R_1\nu(0)) \geq 0$ we get that for all $x \in \Omega$
    \begin{equation*}
        u(x,t) \leq \frac{B\nu(0)}{R_1 \nu(0) + 
        (B - R_1 \nu(0)) e^{B t}} \rightarrow 0 \quad \text{as} \quad t \rightarrow \infty.
    \end{equation*}
    Now, let us replace the quadratic term in equation \eqref{eq:v_stationary} with zero to get
    \begin{equation}
        \label{eq:v_elliptic_linear_part}
        d_v \Delta V - V + A = 0, \quad x \in \Omega, \quad V = 0, \quad x \in \partial \Omega.
    \end{equation}
    Recall that Lemma \ref{lem:existence_of_inhomogeneous_elliptic_v} yields existence of non-constant solution $V_0$ of the problem \eqref{eq:v_elliptic_linear_part}. Denote $w = v - V_0$. One can easily verify that it satisfies the following partial differential equation.
    \begin{equation}
        \label{eq:w_parabolic}
        w_t = d_v\Delta w - w - u^2v \quad x \in \Omega, \quad w = 0 \quad x \in \partial\Omega
    \end{equation}
    Denote by $e^{t(\Delta - Id)}$ the well known heat semigroup. We then apply a Duhamel principle to the equation \eqref{eq:w_parabolic} to obtain
    \begin{equation}
        \label{eq:w_duhmel_principle}
        w(t) = e^{t(d_v\Delta - Id)} w(0) - \int_0^t e^{(t-s)(d_v\Delta - Id)} u^2(s) v(s)\:ds
    \end{equation}
    Equation \eqref{eq:w_duhmel_principle} allows us to estimate the $L^\infty$ norm of $w$ as follows
    \begin{align}
        \label{eq:w_infty_estimate}
        \begin{split}
            \|w(t)\|_\infty 
            &\leq e^{-t}\|w(0)\|_\infty + \int_0^t e^{-(t-s)} \|u^2(s)\|_\infty \|v(s)\|_\infty\:ds \\
            &\leq e^{-t}\|w(0)\|_\infty + R_1\int_0^t e^{-(t-s)} \|u^2(s)\|_\infty \:ds.
        \end{split}
    \end{align}
    Now, we use inequality \eqref{EQ:AprioriEsitmate_Unvariant_Region} to estimate the integrated term as follows
    \begin{align}
        \label{eq:w_infty_estimate_integral}
        \begin{split}
            R_1 \int_0^t e^{-(t-s)} \|u^2(s)\|_\infty \:ds 
            &\leq
            \frac{R_1 B \nu(0)}{B-R_1 \nu(0)} \int_0^t e^{-(t-s)} e^{-Bs}\:ds \\
            &= 
            \frac{R_1 B \nu(0)}{(B-R_1\nu(0))(B-1)} e^{-t} \\
            &- 
            \frac{R_1 B \nu(0)}{(B-R_1\nu(0))(B-1)} e^{-Bt}.
        \end{split}
    \end{align}
    Now, combining inequalities \eqref{eq:w_infty_estimate} and \eqref{eq:w_infty_estimate_integral} we get that
    \begin{align*}
        \|w(t)\|_\infty 
        &\leq \left( \|w(0)\|_\infty + \frac{R_1 B \nu(0)}{(B-R_1\nu(0))(B-1)} \right) e^{-t} \\
        &- \frac{R_1 B \nu(0)}{(B-R_1\nu(0))(B-1)} e^{-Bt}.
    \end{align*}
    It is an immediate consequence that $\|w(t)\|_\infty \rightarrow 0$ as $t \rightarrow \infty$.
\end{proof}

\subsection{De-desertification criteria}

Now, we closely investigate the linearized stability of stationary solutions obtained in Theorem \ref{thm:Stationary-Solutions-Existence}. We start by computing the constant solutions of equations \eqref{eq:u_stationary} and \eqref{eq:v_stationary}. We search for the numbers $(u_*, v_*)$ satisfying
\begin{align}
    \label{EQ:KineticSystem}
    \begin{cases}
        0 = u^2v - Bu,\\
        0 = -u^2v - v + A.
    \end{cases}
\end{align}
\begin{prp}
    \label{prp:ExistenceOfConstantSSKineticSystem}
    Consider parameters $A > 0$ and $B > 0$ of the Klausmeier system \eqref{eq:u}-\eqref{eq:v_init}. The following statements hold true.
    \begin{itemize}
        \item If $A > 2B$ then there are three real solutions of \eqref{EQ:KineticSystem}.
        \item If $A = 2B$ then there are two real solutions of \eqref{EQ:KineticSystem}.
        \item If $A < 2B$ then there is one real solution of \eqref{EQ:KineticSystem}.
    \end{itemize}
\end{prp}

\begin{proof}
    From the second equation we get $v = \frac{A}{u^2 + 1}$. Thus the first equation becomes
    \begin{equation}
        \label{EQ:StationaryPolynomial}
        -B u^3 + A u^2 - B u = 0
    \end{equation}
    For all parameters $A > 0, B > 0$ we can find the solution of following form
    \begin{align}
        \label{EQ:FirstSS}
        \left(u_{*,1}, v_{*,1}\right) &= \left(0, A\right),
    \end{align}
    Provided that $A > 2B$, the equation \eqref{EQ:StationaryPolynomial} admits two more solutions
    \begin{align}
        \left(u_{*,2}, v_{*,2}\right) &= \left(\frac{A - \sqrt{A^2 - 4B^2}}{2B}, \frac{2 B^2}{A - \sqrt{A^2 - 4B^2}}\right), \label{EQ:SecondSS}\\
        \left(u_{*,3}, v_{*,3}\right) &= \left(\frac{A + \sqrt{A^2 - 4B^2}}{2B}, \frac{2 B^2}{A + \sqrt{A^2 - 4B^2}}\right) \label{EQ:ThirdSS}.
    \end{align}
    Finally, if $A = 2B$, then the solutions \eqref{EQ:SecondSS} and \eqref{EQ:ThirdSS} merge into 
    \begin{align}
        \label{EQ:FourthSS}
        \left( u_{*,4}, v_{*,4} \right) &= \left(1, B\right).
    \end{align}
    For $A < 2B$ there are no other, real solutions of \eqref{EQ:StationaryPolynomial} then \eqref{EQ:FirstSS}.
\end{proof}

\begin{rem}
    \label{rem:classification_steady_states}
    The stationary solution $(u_{*,1},v_{*,1})$ corresponds to the desert state, particularly $u_{*,1} = 0$. On the other hand, solutions $(u_{*,2},v_{*,2})$ and $(u_{*,3},v_{*,3})$ are ecologically more interesting as they correspond to equilibria with non-zero biomass.
\end{rem}

Next, we identify the regimes of parameters $A, B$ yielding existence of stable equilibria.

\begin{prp}
    \label{prp:StabilityOfConstanstKineticSystem}
    Denote the following Jacobian matrix of system \eqref{EQ:KineticSystem}
    \begin{equation}
        \label{EQ:JacobianOfConstantSystem}
        D(u, v) = \begin{pmatrix}
            2uv - B & u^2 \\ -2uv & -1 - u^2
        \end{pmatrix}.
    \end{equation}
    Let $i\in \{1,2,3,4\}$ and denote by $(u_{*,i}, v_{*, i})$ the real solution of the system \eqref{EQ:KineticSystem} obtained in Proposition \ref{prp:ExistenceOfConstantSSKineticSystem}. The following statements holds true
    \begin{itemize}
        \item For all $A >0, B > 0$ matrix $D(u_{*,1}, v_{*, 1})$ has only strictly negative eigenvalues.
        \item For all $A > 2 B$ such that $B < 2$ matrix $D(u_{*,3}, v_{*, 3})$ has only strictly negative eigenvalues.
        \item There exists range of parameters $A > 2 B$ such that $B > 2$ and matrix $D(u_{*,3}, v_{*, 3})$ has only strictly negative eigenvalues.
    \end{itemize}
\end{prp}
\begin{proof}
    Substituting  $(u_{*,1}, v_{*,1})$, given by formula \eqref{EQ:FirstSS}, into \eqref{EQ:JacobianOfConstantSystem} we obtain
    \begin{equation*}
        D(u_{*,1}, v_{*,1}) = \begin{pmatrix}
            - B & 0 \\ 0 & -1
        \end{pmatrix} \quad \text{with} \quad \det D(u_{*,1}, v_{*,1}) > 0, \Tr D(u_{*,1}, v_{*,1}) < 0.
    \end{equation*}
    This means that $D(u_{*,1}, v_{*,1})$ has two strictly negative eigenvalues for all $A > 0, B > 0$. Next, observe that 
    \begin{equation*}
        u_{*,i} v_{*,i} = B \quad i \in \{2, 3, 4\}
    \end{equation*}
    It follows that the Jacobian $D(u_{*,i}, v_{*,i})$ given by formula \eqref{EQ:JacobianOfConstantSystem} becomes
    \begin{equation}
        \label{EQ:JacobianOfKineticSystemPlugged}
        D(u_{*,i}, v_{*,i}) = \begin{pmatrix}
            B & (u^*_i)^2 \\ -2B & -1 - (u^*_i)^2
        \end{pmatrix},
        \quad i \in \{2, 3, 4\}
    \end{equation}
    The determinant is given by
    \begin{equation}
        \label{EQ:DeterminantLinearizationKinetic}
        \det D(u_{*,i}, v_{*,i}) = -B -Bu_{*,i}^2 + 2Bu_{*,i}^2 = -B(1-u_{*, i}^2)
    \end{equation}
    Observe that it is positive if and only if $u_i > 1$. For sake of exposition we list if this is satisfied with respect to $i \in \{2, 3, 4\}$.
    \begin{itemize}
        \item $u_{*,2} \leq 1$ for all positive $B$, thus determinant is negative. This yields that one of the eigenvalues of $D(u_{*,2}, v_{*,2})$ must be positive.
        \item $u_{*,3} > 1$ as we assumed $A > 2B$. Hence, we have that the determinant of $D(u_{*,3}, v_{*,3})$ is positive.
        \item $u_{*,4} = 1$ and so the determinant vanishes. This means that one of the eigenvalues of $D(u_{*,4}, v_{*,4})$ is zero.
    \end{itemize}
    Now, the trace of \eqref{EQ:JacobianOfKineticSystemPlugged} can be computed as
    \begin{equation}
        \label{EQ:TraceLinearizationKinetic}
        \Tr D(u_{*,i}, v_{*,i}) = B - 1 - u_{*,i}^2,
    \end{equation}
    which is negative if and only if $(u_i^*)^2 > B - 1$. As we have shown earlier, only $u_3^*$ could be asymptotically stable solution. Observe that $u_3^* > 1$ and consequently if $B < 2$ then the trace is negative, and thus both eigenvalues must be strictly negative real numbers. This means that $u_3^*$ is asymptotically stable. When $B > 2$ we need a stronger lower bound. This bound can be achived for $A > 2$, however itself is nasty, a thus we do not present it.
\end{proof}

\begin{thm}
    \label{thm:Stationary-Solutions-Stability}
    Assume that $d_u > 0$ is sufficiently large and $d_v > 0$ is sufficiently small. Then the following statements hold true
    \begin{itemize}
        \item For all $A > 0, B >0$ the steady state $(U, V)$ constructed around desert state solution $(u_{*,1}, v_{*,1})$ is linearly, asymptotically stable in the sense of the norm of $L^2(\Omega)$.
        \item For all $A > 2B$ such that $B < 2$ the steady state $(U, V)$ constructed around the non-zero solution $(u_{*,3}, v_{*,3})$ is linearly, asymptotically stable in the sense of the norm of $L^2(\Omega)$.
        \item There exists a range of parameters $A > 2B$ such that $B > 2$ and the steady state $(U, V)$ constructed around non-zero solution $(u_{*,3}, v_{*,3})$ is linearly, asymptotically stable in the sense of the norm of $L^2(\Omega)$.
    \end{itemize}
\end{thm}

\begin{proof}
    It is convenient to introduce the following notation
    \begin{equation}
        \label{eq:non_linearities}
        f(u, v) = u^2v - Bu, \quad g(u,v) = -u^2v - v + A
    \end{equation}
    Let $(U, V)$ be a solution constructed in Theorem \ref{thm:Stationary-Solutions-Existence}, around any constant numbers $(u_*, v_*)$ obtained by Proposition \ref{prp:ExistenceOfConstantSSKineticSystem}. We linearize problem \eqref{eq:u}-\eqref{eq:v_bound} around $(U, V)$ and introduce the new functions $\varphi = u - U$ and $\psi = v - V$. This functions satisfy the following system of PDEs.
    \begin{align*}
        \begin{pmatrix}
            \varphi \\ \psi
        \end{pmatrix}_t
        =
        \begin{pmatrix}
            d_u \mathcal{L} & 0 \\ 0 & d_v \Delta
        \end{pmatrix}
        \begin{pmatrix}
            \varphi \\ \psi
        \end{pmatrix}
        +
        \begin{pmatrix}
            f_u(U,V) & f_v(U,V) \\ g_u(U,V) & g_v(U,V)
        \end{pmatrix}
        \begin{pmatrix}
            \varphi \\ \psi
        \end{pmatrix},
    \end{align*}
    where the second matrix is equivalent to the matrix $D(U, V)$, defined by the formula \eqref{EQ:JacobianOfConstantSystem}. From that, we write the energy equation as follows
    \begin{align}
        \label{eq:energy_varphi}
        \frac{\partial}{\partial t} \int_\Omega \varphi^2\:dx 
        &=
        d_u \langle \mathcal{L}\varphi, \varphi \rangle_{L^2(\Omega)} +  \int_\Omega f_u(U,V) \varphi^2 \:dx + \int_\Omega f_v(U,V) \varphi \psi \:dx.
    \end{align}
    Next, we estimate the first term in \eqref{eq:energy_varphi} from above
    \begin{equation}
        \label{eq:energy_varphi_1st_term}
        \langle \mathcal{L}\varphi, \varphi \rangle_{L^2(\Omega)} 
        =
        \frac{\langle \mathcal{L}\varphi, \varphi \rangle_{L^2(\Omega)}}{\|\varphi\|_2^2}\|\varphi\|_2^2
        \leq
        \inf_{\varphi \in L^2(\Omega)}\left( \frac{\langle \mathcal{L}\varphi, \varphi \rangle_{L^2(\Omega)}}{\|\varphi\|^2_2} \right) \|\varphi\|_2^2 =: \beta_1 \|\varphi\|_2^2,
    \end{equation}
    where we denoted by $\beta_1 < 0$ a principle eigenvalue of $\mathcal{L}$ with non-local Dirichlet boundary condition, for details see \cite{GarciaMelianRossi2009}. Furthermore, we deal with the third term in equation \eqref{eq:energy_varphi}, using the well known Cauchy inequality $ab \leq \frac{a^2}{2} + \frac{b^2}{2}$ for arbitrary numbers $a, b \in \mathbb{R}$.
    \begin{equation}
        \label{eq:energy_varphi_2nd_term}
        \int_\Omega f_v(U,V) \varphi \psi \:dx \leq \|f_v(U,V)\|_\infty \left( \frac{\|\varphi\|_2^2}{2} + \frac{\|\psi\|_2^2}{2} \right)
    \end{equation}
    Next, we plug estimates \eqref{eq:energy_varphi_1st_term} and \eqref{eq:energy_varphi_2nd_term} into equation \eqref{eq:energy_varphi} to get that
    \begin{align}
        \label{eq:energy_varphi_odi}
        \begin{split}
            \frac{\partial}{\partial t} \int_\Omega \varphi^2\:dx
            &\leq
            \left(d_u \beta_1 + \frac{\|f_v(U,V)\|_\infty}{2}\right) \|\varphi\|_2^2 \\
            &+
            \int_\Omega f_u(U,V) \varphi^2\:dx + 
            \frac{\|f_v(U,V)\|_\infty}{2}\|\psi\|_2^2.
        \end{split}
    \end{align}
    Analogously, for the energy of $\psi$ we obtain the differential inequality
    \begin{align}
        \label{eq:energy_psi_odi}
        \begin{split}
            \frac{\partial}{\partial t} \int_\Omega \psi^2\:dx
            &\leq
            \left(d_v \lambda_1 + \frac{\|g_u(U,V)\|_\infty}{2} \right) \|\psi\|_2^2 \\
            &+
            \int_\Omega g_v(U,V) \psi^2 \:dx + \frac{\|g_u(U,V)\|_\infty}{2} \|\varphi\|_2^2,
        \end{split}
    \end{align}
    where we denoted by $\lambda_1 < 0$ a principle eigenvalue of $\Delta$ with Dirichlet boundary condition. Combining inequalities \eqref{eq:energy_varphi_odi} and \eqref{eq:energy_psi_odi} we obtain
    \begin{align}
        \label{eq:energy_combined}
        \begin{split}
            \frac{\partial}{\partial t} \int_\Omega \varphi^2 + \psi^2\:dx
            &\leq
            \left(d_u \beta_1 + \frac{\|f_v(U,V)\|_\infty}{2} + \frac{\|g_u(U,V)\|_\infty}{2}\right) \|\varphi\|_2^2 \\
            &+
            \left(d_v \lambda_1 + \frac{\|f_v(U,V)\|_\infty}{2} + \frac{\|g_u(U,V)\|_\infty}{2} \right) \|\psi\|_2^2 \\
            &+ 
            \int_\Omega f_u(U,V) \varphi^2\:dx
            +
            \int_\Omega g_v(U,V) \psi^2 \:dx
        \end{split}
    \end{align}
    From inequality \eqref{eq:energy_combined} it is clear that there is a multivariate relation depending on the magnitude of $d_v$, that yields linear, asymptotic stability of $(U, V)$. In particular, there always exists $d_u, d_v$ large enough, to have that
    \begin{equation}
        \int_\Omega \varphi^2\:dx + \int_\Omega \psi^2\:dx \rightarrow 0, \quad \text{ as } \quad t \rightarrow \infty.
    \end{equation}
\end{proof}
\section{Numerical analysis}

\subsection{Approximation with finite dimensional kernels}
Now, it is convenient to denote the integral kernel $K: \Omega \times \Omega \rightarrow \mathbb{R}_{\geq 0}$ as
\begin{equation}
    \label{EQ:KernelKDef}
    K(x, y) = J(x-y),
\end{equation}
where $J$ is an integral kernel as in Assumption \ref{assumptions_general_kernel}. Observe that $K \in L^2(\Omega \times \Omega)$. Since $L^2(\Omega \times \Omega) = L^2(\Omega) \otimes L^2(\Omega)$ we can expand \eqref{EQ:KernelKDef} into
\begin{equation*}
    K(x, y) = \sum_{k=0}^\infty \sum_{l=0}^\infty c_{k, l} \varphi_k(x) \varphi_l(y),
\end{equation*}
where $\{\varphi_k\}_{k=0}^\infty$ is an orthonormal basis of $L^2(\Omega)$ and $c_{k,l}$ are given coefficients. Hence, we may define the truncated kernel $K_N : \Omega \times \Omega \rightarrow \mathbb{R}_{\geq 0}$
\begin{equation}
    \label{EQ:TruncatedKernelKDef}
    K_N(x, y) = \sum_{k=0}^N \sum_{l=0}^N c_{k, l} \varphi_k(x) \varphi_l(y).
\end{equation}
Analogously we denote the corresponding, truncated dispersal operator $\mathcal{L}_N: C(\Omega) \rightarrow C(\Omega)$ by the formula below
\begin{equation}
    \label{EQ:TruncatedDispersalOperatorDef}
    \mathcal{L}_N u(x, t) = \int_\Omega K_N(x,y) \left(u(y,t) - u(x,t) \right)\:dy
\end{equation}
Finally, we want to study the resulting truncated Klausmeier model
\begin{align}
    \label{EQ:TruncatedKlausmeierSystem}
    \begin{cases}
        u_t(x,t) = d_u\mathcal{L}_N u(x, t) + f(u(x,t), v(x,t)) &(x, t) \in \Omega \times [0, T] \\
        v_t(x,t) = d_v\Delta v(x,t) + g(u(x,t), v(x,t)) &(x, t) \in \Omega \times [0, T] \\
        v(x,t) = \Bar{v} \in \mathbb{R}_{\geq 0} &(x, t) \in \partial\Omega \times [0, T]. \\
        u(x,0) = u_0(x), v(x,0) = v_0(x) &x \in \Omega,
    \end{cases}
\end{align}
where $f, g, d_v, \nu$ are the same as in \eqref{EQ:KlausmeierSystem}. Now, we shall prove the convergence of system \eqref{EQ:TruncatedKlausmeierSystem} to \eqref{EQ:KlausmeierSystem} as $N \rightarrow \infty$.
\begin{thm}
    \label{thm:ApproximationWithFiniteKernels}
    Consider the initial data 
    $$(u_0, v_0) \in C(\Omega) \times \{v \in H^1_0(\Omega) \cap C(\Omega): \Delta v \in C(\Omega)\}$$
    Let us denote the solution of the problem \eqref{EQ:TruncatedKlausmeierSystem} with maximal time of evolution $T_N > 0$ such that
    $$(u_N, v_N) \in C^{0, 1}(\Omega, [0,T]) \times C^{0, 1}(\Omega, [0,T])$$
    Let us also denote the solution of the problem \eqref{EQ:KlausmeierSystem} with maximal time of evolution $T > 0$ such that
    $$(u, v) \in C^{0, 1}(\Omega, [0,T]) \times C^{0, 1}(\Omega, [0,T])$$
    Then
    \begin{equation*}
        (u_N, v_N) \xrightarrow{N \rightarrow \infty} (u, v),
    \end{equation*}
    where the convergence is understood in the sense of the product norm $||| \cdot |||$ defined by the formula
    \begin{equation}
        ||| (u, v) ||| = \sup_{t \in [0, T_{min}]}\left(\| u(\cdot, t) \|_2 + \|v(\cdot, t) \|_2 \right).
    \end{equation}
\end{thm}
\begin{proof}
    Let us denote $U_N = u - u_N$ and $V_n = v - v_n$. Next, we plug $U_N$ and $V_N$ into \eqref{EQ:KlausmeierSystem}. Observe that $U_N$ satisfies the following equation
    \begin{align*}
        \partial_t U_N 
        &= d_u \mathcal{L}_N U_N \\
        &+ d_u \sum_{k > N}\sum_{l > N} c_{k, l}\int_\Omega \varphi_k(x, s) \varphi_l(y, s) (u(y, s) - u(x, s))\:dy \\
        &+ f(u(x,s), v(x,s)) - f(u_N(x,s), v_N(x,s)).
    \end{align*}
    Multiplying this equation by $U_N$ and integrating in $\Omega$, we obtain the inequality
    \begin{align}
        \label{EQ:U_NDifferentialInequality_PRE}
        \begin{split}
            \frac{1}{2}\partial_t\|U_N(\cdot, t)\|_2^2 
            &\leq 2 b_\infty d_u  \langle \mathcal{L} U_N(\cdot, s), U_N(\cdot, s)\rangle \\
            &+ d_u |\Omega| \max_{t \in [0, T]}\|u(\cdot, t)\|_2^2 \sum_{k > N}\sum_{l > N} c_{k, l} \\
            &+ \|f_u\|_\infty \|U_N(\cdot, s)\|_2^2  + \|f_v\|_\infty \|V_N(\cdot, s)\|_2^2.
        \end{split}
    \end{align}
    Observe that $\mathcal{L}$ annihilates constants and is a conservative operator. Thus, we estimate the first term by
    \begin{equation*}
        \langle \mathcal{L} U_N(\cdot, s), U_N(\cdot, s)\rangle
        \leq
        \sup_{\|w\|_2 = 1, \int_\Omega w \:dx = 0} \langle \mathcal{L}w, w \rangle \|U_N(\cdot, s)\|_2^2 = \beta_1 \|U_N(\cdot, s)\|_2^2,
    \end{equation*}
    where $\beta_1 < 0$, see \cite{ChasseigneChavesRossi2006} for details. Plugging this into \eqref{EQ:U_NDifferentialInequality_PRE} we obtain that
    \begin{align}
        \label{EQ:U_NDifferentialInequality_Final}
        \begin{split}
            \frac{1}{2}\partial_t\|U_N(\cdot, t)\|_2^2
            &\leq (2 b_\infty d_u \beta_1 + \|f_u\|_\infty) \|U_N(\cdot, t)\|_2^2 + \|f_v\|_\infty \|V_N(\cdot, t)\|_2^2 \\
            &+ d_u |\Omega| \max_{t \in [0, T]}\|u(\cdot, t)\|_2^2 \sum_{k > N}\sum_{l > N} c_{k, l}.
        \end{split}
    \end{align}
    Proceeding similarly with $V_N$ one finds out that
    \begin{equation}
        \label{EQ:V_NDifferentialInequality_Final}
            \frac{1}{2}\partial_t\|V_N(\cdot, t)\|_2^2 \leq ( d_v \lambda_1 + \|g_v\| ) \|V_N(\cdot, t)\|_2^2 + \|g_u\|_\infty \|U_N(\cdot, t)\|_2^2.
    \end{equation}
    Now, combining inequalities \eqref{EQ:U_NDifferentialInequality_Final} and \eqref{EQ:V_NDifferentialInequality_Final} we obtain further bounds
    \begin{align}
        \label{EQ:U_NV_N_FinalDifferentialInequality}
        \begin{split}
            \frac{1}{2}\partial_t( \|V_N(\cdot, t)\|_2^2 + \|U_N(\cdot, t)\|_2^2 )
            &\leq M ( \|V_N(\cdot, t)\|_2^2 + \|U_N(\cdot, t)\|_2^2 ) \\
            &+ d_u |\Omega| \max_{t \in [0, T]}\|u(\cdot, t)\|_2^2 \sum_{k > N}\sum_{l > N} c_{k, l}.
        \end{split}
    \end{align}
    where we denoted the following finite, real number 
    $$M = \max\left\{ 2 b_\infty d_u \beta_1 + \|f_u\|_\infty + \|g_u\|_\infty,\: d_v \lambda_1 + \|g_v\|_\infty + \|f_v\|_\infty \right\}$$
    Now, integrating \eqref{EQ:U_NV_N_FinalDifferentialInequality} in time and applying the integral Gr\"onwall inequality we obtain that
    \begin{equation*}
        ||| (U_N, V_N) ||| \leq d_u |\Omega| T_{min} \max_{t \in [0, T_{min}]}\|u(\cdot, t)\|_2^2 e^{M t} \sum_{k > N}\sum_{l > N} c_{k, l}.
    \end{equation*}
    We end the proof of this Theorem by passing to the limit with $N \rightarrow \infty$.
\end{proof}

\subsection{Numerical scheme}

Here, we consider $\{ P_k \}_{k \geq 0}$ to be the set of Legendre polynomials. Note that this is an orthonormal set on $L^2(\Omega)$. Since they are dense in $C(\Omega)$ we can approximate $K$ as tensorial product of Legendre polynomials, e.g. there exists sequence of symmetric coefficients $c_{k,l} = c_{l, k}$ such that
\begin{equation}
    K(x, y) = \sum_{k \geq 0} \sum_{l \geq 0} c_{k, l} P_k(x) P_l(y)
\end{equation}
Theorem \ref{thm:ApproximationWithFiniteKernels} yields that for any $\varepsilon > 0$ there exists a natural number $0 < N < \infty$ such that the solution $(u_N, v_N)$ of the problem \eqref{EQ:TruncatedKlausmeierSystem} and solution $(u, v)$ of the system \eqref{EQ:KlausmeierSystem} satisfy
\begin{equation}
    |||(u_N, v_N) - (u, v)||| \leq \varepsilon.
\end{equation}
Thus we will numerically solve the truncated system \eqref{EQ:TruncatedKlausmeierSystem}. We particularly look for solutions of the following form
\begin{equation}
    \label{EQ:UDefinedAsLegendrePolys}
    u(x, t) = \sum_{i, j = 0}^N a_{i, j} P_i(x) P_j(t).
\end{equation}
Recall that the dispersal operator decomposes into sum of integral operator and multiplication operator as follows
\begin{equation}
    \mathcal{L}_N f(x) = \int_\Omega K_N(x,y) f(y)\:dy - b(x) f(x).
\end{equation}
Now, let us plug the formula \eqref{EQ:UDefinedAsLegendrePolys} into the first term
\begin{align*}
    \int_\Omega K_N(x,y) u(x,t) \:dy
    &= \sum_{i, j, k, l = 0}^N a_{i, j} c_{k, l} P_k(x) P_j(t) \int_\Omega P_l(y) P_i(y) \:dy \\
    &= \sum_{j, k = 0}^N \Tilde{a}_{j, k} P_k(x) P_j(t),
\end{align*}
where we denoted the coefficients
\begin{equation*}
    \Tilde{a}_{j ,k} = \sum_{i=1}^N a_{i, j} c_{k, i}
\end{equation*}
\subsection{Stability of the numerical scheme}

\subsection{Experiments}

\newpage

\bibliographystyle{siam}
\bibliography{bibliography}

% \begin{thebibliography}{100}

% \end{thebibliography}

\end{document}